\section{Kernel implementations}

The suite is a ladder along the arithmetic intensity axis. Each rung is a kernel
whose position on the roofline is meant to teach something specific.

\subsection{SAXPY}

$y \leftarrow a x + y$ over long vectors. Two floating point operations per
element against three memory accesses (two reads, one write), so its intensity is
tiny and it lives at the far memory bound left edge of the roofline. It is the
anchor: if SAXPY does not reach close to the measured bandwidth ceiling, the
harness or the copy microbenchmark is wrong before any GEMM is even considered.

\subsection{Reduction}

A shared memory tree reduction that sums a large array. It exercises the
synchronization and shared memory path rather than raw streaming, and it is
tested across non power of two sizes because the tail handling is where reduction
kernels usually break.

\subsection{Transpose, naive and tiled}

The transpose does no floating point work, which isolates coalescing and shared
memory bank conflicts with nothing else in the way. The naive version reads the
input coalesced but writes the transpose with a large stride, so the writes are
uncoalesced. The tiled version stages a tile in shared memory so both the global
read and the global write are coalesced. The tile is declared one column wider
than it is tall, \texttt{tile[TILE\_DIM][TILE\_DIM+1]}, so that the transposed
access pattern to shared memory does not map every thread in a warp onto the same
bank. That single pad word removes the bank conflicts on the transposed access,
and the naive versus tiled delta is the cleanest demonstration in the whole
suite.

\subsection{GEMV}

Matrix times vector. Each matrix element is read once and touched by one
multiply add, so reuse is low and the kernel sits in the middle of the intensity
axis, between the streaming kernels and the GEMM ladder.

\subsection{The GEMM ladder}

General matrix multiply is where intensity and achieved performance climb
together as reuse improves, and it is the core of the report.

\begin{description}
  \item[Naive.] Every thread reads its row and column straight from global
    memory. Correct, and slow, because each value is reloaded from DRAM many
    times. This is the memory bound bottom of the GEMM ladder.
  \item[Shared memory tiled.] Threads cooperate to stage tiles of $A$ and $B$
    into shared memory once, then reuse them across the tile. Reuse rises and the
    kernel climbs the roof.
  \item[Register blocked.] Each thread computes a small output tile rather than a
    single element, holding partial sums in registers. This raises the reuse per
    shared memory load, trading register pressure for far fewer shared accesses,
    and it is usually the step with the largest jump in achieved performance.
  \item[Vectorized.] The register blocked kernel with \texttt{float4} loads and
    stores, so each memory instruction moves four values and the load and store
    issue pressure drops.
  \item[cuBLAS SGEMM.] The vendor library, plotted as an honest ceiling for a
    custom kernel to be measured against. It is labeled as a library, not as one
    of my kernels, because it is not a fair comparison of hand written code and
    it would be dishonest to present it as one.
\end{description}

A WMMA tensor core GEMM is a stretch goal on its own roofline panel, since its
compute ceiling is a different number from the FP32 CUDA core ceiling and the two
should not share an axis.
